\chapter{Evaluation}
\label{ch:evaluation}
Dieses Kapitel präsentiert die Durchführung (\autoref{subsec:messbedingungen}) sowie die Ergebnisse der Metriken (\autoref{subsec:metriken}), um zu beurteilen, inwiefern sich die lokalen Modelle für die Verarbeitung von EEG-Daten eignen.

\section{Durchführung der Messungen}
\label{sec:eval-durchfuehrung}
Wie schon in \autoref{sec:messprotokoll} beschrieben, wird die Eingabe für alle Modelle identisch gehalten. Insgesamt werden so alle neun Kombinationen aus drei Modellen und drei Szenarien auf dem Samsung Galaxy S23 vermessen und anschließend in einer neu erstellten CSV-Datei protokolliert. Dies geschieht je fünfmal, wobei die erste Messung jeder Kombination verworfen wird.
Vor der ersten Messung wird ein Aufwärmlauf durchgeführt.
\medskip

\noindent Als Eingabe für das Modell werden zwei bestimmte Sessions verwendet. Beide sind beim Tragen des Unicorn-Hybrid-Headsets entstanden und haben eine Laufzeit von knapp über 120 Sekunden.
% Kurzer Recap (kein Wiederholen der ganzen Methodik):
% - 3 Modelle x 3 Szenarien, N=5, erste Messung verworfen
% - Bedingungen -> \autoref{subsec:messbedingungen}
% - Testgerat S23
% 3-4 Satze reichen.

\section{Latenz}
\label{sec:eval-latenz}
Im folgenden werden die Ergebnisse der Latenzmessungen vorgestellt.
% Deutung: welches Modell schneller? TTFT vs Gesamtdauer? Tokens/s als fairer Vergleich.
\setlength{\tabcolsep}{4pt}
\begin{table}[H]
    \centering
    \caption{Latenzmetriken je Modell und Szenario (Mittel aus vier Wiederholungen $\pm$ Standardabweichung).}
    \label{tab:eval-latenz}
    \begin{tabular}{llrrrr}
        \toprule
        \textbf{Modell} & \textbf{Szenario} & \textbf{Ladezeit (s)} & \textbf{TTFT (s)} & \textbf{Gesamtdauer (s)} & \textbf{Tok/s} \\
        \midrule
        \multirow{3}{*}{Gemma 3 1B}
            & Report    & 0,45 & 1,94 $\pm$ 0,07 & 6,70 $\pm$ 0,14  & 21,65 $\pm$ 0,44 \\
            & Artefakt  & 0,34 & 1,95 $\pm$ 0,04 & 11,44 $\pm$ 0,04 & 25,34 $\pm$ 0,09 \\
            & Vergleich & 0,29 & 5,86 $\pm$ 0,49 & 9,89 $\pm$ 0,46  & 11,85 $\pm$ 0,54 \\
        \midrule
        \multirow{3}{*}{Qwen 2.5 1.5B}
            & Report    & 0,73 & 5,36 $\pm$ 0,10 & 34,43 $\pm$ 0,24 & 14,06 $\pm$ 0,10 \\
            & Artefakt  & 0,27 & 5,31 $\pm$ 0,13 & 26,70 $\pm$ 0,12 & 13,37 $\pm$ 0,06 \\
            & Vergleich & 0,27 & 13,99 $\pm$ 0,55 & 42,00 $\pm$ 1,94 & 11,31 $\pm$ 0,53 \\
        \midrule
        \multirow{3}{*}{Gemma 2 2B}
            & Report    & 0,87 & 6,09 $\pm$ 0,49 & 14,25 $\pm$ 0,49 & 5,06 $\pm$ 0,17 \\
            & Artefakt  & 0,47 & 6,75 $\pm$ 0,38 & 18,57 $\pm$ 0,58 & 4,74 $\pm$ 0,15 \\
            & Vergleich & 0,52 & 16,60 $\pm$ 0,23 & 27,19 $\pm$ 0,13 & 2,87 $\pm$ 0,01 \\
        \bottomrule
    \end{tabular}
\end{table}

\noindent \autoref{tab:eval-latenz} zeigt, dass der \textbf{Durchsatz (Tok/s)} bei dem Modell Gemma 3 1B am höchsten ist (bis zu 25,34 Tok/s), während das Modell Gemma 2 2B den geringsten Durchsatz aufweist (minimal 2,87 Tok/s).
Die \textbf{Ladezeit} ist bei Gemma 2 2B am höchsten; die beiden kleineren Modelle Gemma 3 1B und Qwen 2.5 1.5B laden schneller. Sie liegt durchweg unter einer Sekunde (zwischen 0,27 und 0,87,s).
Die \textbf{Zeit bis zum ersten Token (TTFT)} ist bei Gemma 3 1B am geringsten, während sie bei Gemma 2 2B am höchsten ist. Mit Ausnahme von Qwen 2.5 1.5B/Artefakt steigt die TTFT mit dem Szenarioniveau an und erreicht ihren Höchstwert bei Gemma 2 2B mit 16,60,s.
Die \textbf{Gesamtdauer} reicht von etwa 6,70 bis 42,00,s und variiert stark zwischen den Modellen und Szenarien.
\medskip

\noindent Da die \textbf{Standardabweichungen} überwiegend gering ausfallen, weisen sie auf stabile Messungen hin. Lediglich bei Qwen 2.5 1.5B/Vergleich ist die Standardabweichung mit 1,94,s (Gesamtdauer) deutlich höher als bei allen anderen Messungen.
\section{Ressourcenverbrauch}
\label{sec:eval-ressourcen}
Daraufolgend werden die Ergebnisse der Messungen zum Ressourcenverbrauch vorgestellt.
% Deutung: RAM waechst mit Modellgroesse? Akku-Tradeoff (2B ~1,5x 1B, aus deinem 1. Lauf).

\begin{table}[H]
    \centering
    \caption{Ressourcenverbrauch je Modell (RAM als Mittel $\pm$ Standardabweichung).}
    \label{tab:eval-ressourcen}
    \begin{tabular}{llrr}
        \toprule
        \textbf{Modell} & \textbf{Szenario} & \textbf{RAM / PSS (MB)} & \textbf{Akku (mAh)\textsuperscript{*}} \\
        \midrule
        \multirow{3}{*}{Gemma 3 1B}
            & Report    & 1482 $\pm$ 43 & \multirow{3}{*}{\textasciitilde70} \\
            & Artefakt  & 1508 $\pm$ 10 & \\
            & Vergleich & 1615 $\pm$ 35 & \\
        \midrule
        \multirow{3}{*}{Qwen 2.5 1.5B}
            & Report    & 1994 $\pm$ 2  & \multirow{3}{*}{\textasciitilde200} \\
            & Artefakt  & 1994 $\pm$ 1  & \\
            & Vergleich & 2114 $\pm$ 34 & \\
        \midrule
        \multirow{3}{*}{Gemma 2 2B}
            & Report    & 3691 $\pm$ 129 & \multirow{3}{*}{\textasciitilde100} \\
            & Artefakt  & 3638 $\pm$ 168 & \\
            & Vergleich & 3798 $\pm$ 128 & \\
        \bottomrule
    \end{tabular}
    \\[4pt]
    \footnotesize\textsuperscript{*}Verbrauch über den gesamten Messblock (15 Läufe) des jeweiligen Modells.
\end{table}

\noindent \autoref{tab:eval-ressourcen} zeigt, dass der \textbf{RAM (PSS)} eine klare Rangfolge zwischen den Modellen besitzt. Während Gemma 3 1B mit (1482 MB) den geringsten Wert aufweist, liegt Qwen in der Mitte (ca.\ 2000 MB) und Gemma 2 2B am höchsten (3798 MB).
Während das 1B-Modell von Gemma mit (ca.\ 70 mAh) den geringsten \textbf{Akkuverbrauch} aufweist, liegt der Verbrauch von Qwen 2.5 1.5B (ca.\ 200 mAh) und Gemma 2 2B (ca.\ 100 mAh) deutlich höher.
Auffällig ist, dass das Modell Qwen am meisten Akku verbraucht.
\section{Stabilität}
\label{sec:eval-stabilitaet}
Alle 45 Messungen wurden mit dem Status \enquote{ok} protokolliert. Das bedeutet, dass alle Läufe erfolgreich liefen. Es gab keine Abstürze, keine Fehler und keinen Speicherüberlauf ().
Das größere Modell Gemma 2 2B lief ebenso stabil, obwohl es mit ca.\ 3742,MB den meisten RAM verbrauchte.
Auffällig war lediglich eine leichte Schwankung, die bei Gemma 2 2B/Artefakt auftrat. Dort fiel eine größere Streuung auf (± 2,24,s), die vermutlich unter Dauerlast entstand.
% Status-Verteilung (erfolgreich/Fehler/OOM), Throttling-Beobachtungen,
% ob ein Modell an die RAM-Grenze des S23 kam. Meist Fliesstext + evtl. kleine Tabelle.

\section{Ausgabequalität}
\label{sec:eval-qualitaet}
% Qualitative Bewertung. Erst Kriterien definieren, dann Beispielausgaben einordnen.

\subsection{Bewertungskriterien}
\label{subsec:eval-kriterien}
\begin{itemize}
    \item \textbf{Korrektheit:} Stimmt die Aussage mit den Merkmalen überein?
    \item \textbf{Verständlichkeit:} Ist die Ausgabe für Laien lesbar?
    \item \textbf{Relevanz:} Geht sie auf das Wesentliche ein?
    \item \textbf{Halluzination:} Werden Dinge erfunden, die nicht in den Daten stehen?
\end{itemize}

\subsection{Beispielausgaben je Szenario}
\label{subsec:eval-beispiele}
% Pro Szenario 1-2 Beispielausgaben (z.B. als \begin{quote} oder lstlisting)
% + kurze Einordnung anhand der Kriterien oben.
% Szenario 1: Session-Report
% Szenario 2: Artefakt-Erklarung
% Szenario 3: Session-Vergleich

\section{Zusammenfassende Gegenüberstellung}
\label{sec:eval-vergleich}
% Ein Gesamtbild: welches Modell fuer welchen Zweck? Tabelle oder Ampel-Uebersicht.
% Ueberleitung zur Diskussion (\autoref{ch:diskussion}).